% ============================================================
\section{Methods}
\label{sec:methods}
% ============================================================

This section describes the complete \kmerdet{} variant detection algorithm from
end to end.  Each subsection first describes the corresponding step in the
original \km{} implementation, then presents the \kmerdet{} enhancements.
Algorithm~\ref{alg:detect} gives a high-level pseudocode summary of the full
pipeline.

\begin{algorithm}[t]
\SetAlgoLined
\KwIn{Target FASTA file $T$; \jf{} database $\mathcal{D}$; configuration $\theta$}
\KwOut{Set of \texttt{VariantCall} records $\mathcal{V}$}

$\mathcal{V} \leftarrow \emptyset$\;
\ForEach{target sequence $t \in T$ \textbf{(in parallel via Rayon)}}{
  $R \leftarrow$ \textsc{DecomposeTarget}$(t, k)$  \tcp*{reference k-mers, in order}
  \If{adaptive thresholds enabled}{
    $\theta_t \leftarrow$ \textsc{AdaptiveThreshold}$(\mathcal{D}, R, \theta)$\;
  }
  \If{bidirectional walking}{
    $W \leftarrow$ \textsc{WalkForward}$(\mathcal{D}, R, \theta_t)$ $\cup$
                   \textsc{WalkBackward}$(\mathcal{D}, R, \theta_t)$\;
  }\Else{
    $W \leftarrow$ \textsc{WalkForward}$(\mathcal{D}, R, \theta_t)$\;
  }
  $W \leftarrow$ \textsc{PruneGraph}$(W)$  \tcp*{tip removal, bubble collapsing}
  $G \leftarrow$ \textsc{BuildGraph}$(W, R)$  \tcp*{directed weighted k-mer graph}
  $\mathcal{P} \leftarrow$ \textsc{FindPaths}$(G)$  \tcp*{Dijkstra; ref edge removal; all-shortest-paths}
  \ForEach{alternative path $p \in \mathcal{P}$}{
    $c \leftarrow$ \textsc{Classify}$(R, p)$  \tcp*{diff\_path\_without\_overlap}
    $(f, \sigma) \leftarrow$ \textsc{QuantifyNNLS}$(W, R, p)$  \tcp*{rVAF + bootstrap CI}
    $q \leftarrow$ \textsc{StatisticalScore}$(W, R, p, \mathcal{D})$  \tcp*{binomial p-value + Phred QUAL}
    $\mathcal{V} \leftarrow \mathcal{V} \cup \{$\textsc{MakeCall}$(t, p, c, f, \sigma, q)\}$\;
  }
}
\Return $\mathcal{V}$\;
\caption{kmerdet detect: full variant detection pipeline}
\label{alg:detect}
\end{algorithm}

\subsection{K-mer Encoding}
\label{sec:kmer-encoding}

All k-mers in \kmerdet{} are stored as 2-bit packed 64-bit integers.  The
encoding assigns A$=00$, C$=01$, G$=10$, T$=11$ to each base, packing the
k-mer into the least-significant $2k$ bits of a \texttt{u64}.  For $k = 31$,
this occupies 62 bits, fitting in a single 64-bit register with 2 bits to spare.

\begin{equation}
  \text{encode}(b_1 b_2 \cdots b_k) = \sum_{i=1}^{k} \text{enc}(b_i) \cdot 4^{k-i}
  \label{eq:kmer-encode}
\end{equation}

The core \texttt{Kmer} type in \kmerdet{} is:

\begin{lstlisting}[style=rust, caption={The core Kmer type (src/kmer/mod.rs).}]
pub struct Kmer {
    pub data: u64,  // 2-bit packed bases, right-aligned
    pub k: u8,      // k-mer length (up to 32)
}

impl Kmer {
    pub fn extend_right(&self, base: u8) -> Self {
        let shifted = (self.data << 2) | (base as u64 & 0b11);
        Self::new(shifted, self.k)   // mask to 2k bits
    }

    pub fn canonical(&self) -> Self {
        // min(self, self.reverse_complement())
        canonical::canonical(*self)
    }
}
\end{lstlisting}

\subsubsection{Canonical Form}

\jf{} in canonical mode stores each k-mer and its reverse complement under the
lexicographically smaller representation.  To match \jf{}'s canonical form
exactly, \kmerdet{} computes the reverse complement by:
\begin{enumerate}
  \item Reversing the bit order of the $2k$ occupied bits (bitwise reversal in 64 bits).
  \item Complementing each 2-bit pair (XOR with $11_2$ per pair, equivalent to XOR with \texttt{0x5555...}).
\end{enumerate}
The canonical form is then $\min(\text{forward}, \text{reverse complement})$ where
comparison is performed on the raw \texttt{u64} value.

\subsection{Target Loading and Reference Sequence Decomposition}
\label{sec:target-loading}

Targets are short FASTA sequences (typically 150--500~bp) centred on the variant
site of interest, with flanking reference sequence that provides unique anchor
k-mers at each end.  The flanking length is typically 35~bp per side (matching
the \kam{} \texttt{multiseqex --flank 35} parameter), which is sufficient to
guarantee a unique $k=31$ anchor for most genomic regions.

For each target sequence, \kmerdet{} decomposes the reference into an ordered
list of $n - k + 1$ overlapping k-mers where $n$ is the target length.  These
reference k-mers serve four roles:
\begin{enumerate}
  \item \textbf{Walking seeds}: DFS extension starts from each reference k-mer.
  \item \textbf{Reference path}: The ordered list encodes the expected (wild-type) sequence.
  \item \textbf{Coverage estimation}: Reference k-mer counts from $\mathcal{D}$ provide local sequencing depth.
  \item \textbf{Graph anchors}: The first and last reference k-mers serve as the virtual source (BigBang) and sink (BigCrunch) of the pathfinding graph.
\end{enumerate}

A target validity check verifies that the anchor k-mers occur uniquely in the
database: if either anchor has count zero, the target is skipped with a
diagnostic warning; if either is non-unique (count far exceeding neighbour
counts without a clear signal), the result is flagged as potentially unreliable.

\subsection{K-mer Walking: Depth-First Extension}
\label{sec:walking}

\subsubsection{Forward Extension}

The walking algorithm discovers all k-mers reachable from the reference by
iterated single-base extension.  For k-mer $S = s_1 s_2 \cdots s_k$, the four
forward extensions are $s_2 s_3 \cdots s_k b$ for $b \in \{A, C, G, T\}$.
Each extension is queried in $\mathcal{D}$, and extensions above the threshold
are added to the DFS stack.  The extension threshold is:

\begin{equation}
  \tau = \max\!\bigl(\lfloor \sigma_{\text{sib}} \cdot r \rfloor,\; n_{\min}\bigr)
  \label{eq:threshold}
\end{equation}

where $\sigma_{\text{sib}}$ is the sum of counts of all four sibling extensions,
$r$ is the ratio threshold (default $r = 0.05$ in \kmerdet{}, $0.00001$ for
maximum sensitivity), and $n_{\min}$ is the absolute count floor (default 2).

The iterative DFS implementation maintains an explicit stack of (k-mer, break\_count)
tuples, bounded by \texttt{max\_stack} entries in the stack (default 500) and
\texttt{max\_break} branch-point crossings (default 10).  A branch point
occurs when $|\text{children}| > 1$.  The total number of discovered nodes is
capped by \texttt{max\_node} (default 10,000).

\subsubsection{Bidirectional Walking}

Forward-only walking can miss variants near the beginning of a target, because
there are insufficient upstream reference k-mers from which to extend.  Similarly,
large deletions that span the right portion of a target may not be discoverable
by walking left-to-right.

\kmerdet{} adds backward extension, which prepends a base: from k-mer $S$,
the four backward extensions are $b \cdot s_1 s_2 \cdots s_{k-1}$ for each
$b$.  Walking both directions and taking the union of discovered nodes improves
sensitivity for boundary variants and is computationally negligible (each
backward walk is bounded identically to forward).

\begin{equation}
  W_{\text{bidir}} = W_{\text{fwd}} \cup W_{\text{bwd}}
  \label{eq:bidir}
\end{equation}

When both directions discover the same k-mer, the maximum count is retained.

\subsubsection{Coverage-Adaptive Thresholds}
\label{sec:adaptive}

The fixed ratio $r$ in equation~\eqref{eq:threshold} fails at coverage extremes.
At $5000\times$ coverage with $0.1\%$ VAF, the reference count is $\sim$5000
and the variant count is $\sim$5.  The ratio $5/5000 = 0.001$, far below the
standard $r = 0.30$ used by \km{} in RNA-seq mode, but the absolute minimum
$n_{\min} = 2$ saves this branch.  At $50000\times$ coverage, error k-mers
have expected counts of $\sim$5 (at $0.1\%$ per-base error rate), so $n_{\min}
= 2$ would pass many error branches.

\kmerdet{}'s adaptive mode estimates the local coverage $C$ from the median
reference k-mer count in a sliding window of width 15, then sets the threshold
from a Poisson noise model~\cite{kelley2010quake,guo2024statistically}:

\begin{equation}
  \lambda_{\varepsilon} = C \cdot \varepsilon / 3
  \qquad
  \tau_{\text{adapt}} = \min\!\bigl\{t : P(\text{Poisson}(\lambda_\varepsilon) \geq t) < \mathrm{FER}\bigr\}
  \label{eq:adaptive}
\end{equation}

where $\varepsilon \approx 0.001$ is the per-base error rate after UMI
deduplication and FER is the desired false extension rate (default $0.01$).
The floor $n_{\min} = 2$ is always enforced for safety.

\subsubsection{Graph Pruning: Tips, Bubbles, and Quick-Rejoin Branches}
\label{sec:pruning}

Before graph construction, three post-walking pruning steps reduce graph complexity:

\textbf{Tip removal}: A ``tip'' is a chain of non-reference k-mers that
terminates without rejoining the reference or any other high-coverage path.
Tips shorter than $\lfloor k/2 \rfloor$ nodes are almost certainly sequencing
errors (a real variant must change at least 1 and at most $k$ consecutive k-mers
for substitutions, more for indels).  They are removed from $W$ before graph
construction.

\textbf{Bubble collapsing}: A ``bubble'' is a pair of paths that share the same
predecessor and successor nodes but differ by one or a few bases at interior
positions.  When one path has count $<10\%$ of the other and the bubble spans
fewer than $k$ nodes, the low-count path is removed.  Bubbles spanning $\geq k$
nodes are preserved because they represent genuine variant paths.

\textbf{Quick-rejoin pruning}: A non-reference extension that rejoins the
reference within fewer than $k$ positions cannot represent a real SNV (which
changes exactly $k$ k-mers) and is a sequencing-error shortcut.  These branches
are removed.

\subsection{Graph Construction}
\label{sec:graph}

From the pruned walk result $W$, \kmerdet{} builds a directed weighted graph
$G = (V, E, w)$:

\begin{itemize}
  \item \textbf{Nodes} $V$: Each discovered k-mer plus virtual BigBang (source)
        and BigCrunch (sink) nodes.
  \item \textbf{Edges} $E$: A directed edge $u \to v$ exists if the $(k-1)$-suffix
        of $u$ equals the $(k-1)$-prefix of $v$.
  \item \textbf{Weights} $w$: Reference edges receive weight $w_{\text{ref}} = 0.01$;
        non-reference edges receive weight $w_{\text{alt}} = 1.0$.  Virtual-node
        edges connect BigBang to each k-mer with no outgoing reference predecessors
        and BigCrunch to each k-mer with no outgoing successors.
\end{itemize}

The asymmetric weighting ($100\times$ lower for reference edges) encodes the
prior that the reference sequence is by far the most likely path through the
graph, ensuring Dijkstra's algorithm recovers the reference path as the globally
shortest path.

\subsection{Pathfinding: Dijkstra, Reference Removal, All-Shortest-Paths}
\label{sec:pathfinding}

The pathfinding procedure follows the original \km{} approach~\cite{audoux2020km}:

\begin{enumerate}
  \item Run Dijkstra from BigBang, recording the predecessor of each node in the
        shortest path tree.
  \item Run Dijkstra from BigCrunch on the reversed graph, recording successors.
  \item Remove all reference edges from $G$.
  \item For each remaining non-reference edge $u \to v$, reconstruct the full
        path from BigBang through $u \to v$ to BigCrunch by concatenating the
        BigBang-to-$u$ prefix and the $v$-to-BigCrunch suffix.
  \item Deduplicate paths; retain each unique sequence.
\end{enumerate}

This procedure efficiently finds all ``shortest deviation paths'' --- the
minimum-weight paths from source to sink that include at least one non-reference
edge.  The two-pass Dijkstra design avoids the combinatorial explosion of
enumerating all paths by pre-computing prefix and suffix distances.

Figure~\ref{fig:graph-example} illustrates the graph for a simple SNV.

\begin{figure}[t]
  \centering
  \begin{tikzpicture}[
    node distance=1.8cm,
    knode/.style={circle, draw=gray!70, fill=gray!10, minimum size=0.9cm, font=\footnotesize},
    refedge/.style={-{Latex[length=2mm]}, gray, thick},
    altedge/.style={-{Latex[length=2mm]}, BrickRed, thick},
    virtnode/.style={rectangle, draw=black!50, fill=black!5, minimum height=0.8cm, font=\footnotesize}
  ]

  \node[virtnode] (bb) {BigBang};
  \node[knode, right=of bb] (k1) {$K_1$};
  \node[knode, right=of k1] (k2) {$K_2$};
  \node[knode, above right=0.9cm and 1.5cm of k2] (kv) {\textcolor{BrickRed}{$K_v$}};
  \node[knode, below right=0.9cm and 1.5cm of k2] (k3ref) {$K_3$};
  \node[knode, right=2.5cm of k2] (k4) {$K_4$};
  \node[virtnode, right=of k4] (bc) {BigCrunch};

  % Reference edges
  \draw[refedge] (bb) -- (k1) node[midway,above,font=\tiny] {0.01};
  \draw[refedge] (k1) -- (k2) node[midway,above,font=\tiny] {0.01};
  \draw[refedge] (k2) -- (k3ref) node[midway,below,font=\tiny] {0.01};
  \draw[refedge] (k3ref) -- (k4) node[midway,below,font=\tiny] {0.01};
  \draw[refedge] (k4) -- (bc) node[midway,above,font=\tiny] {0.01};

  % Variant edge
  \draw[altedge] (k2) -- (kv) node[midway,above,font=\tiny] {1.0};
  \draw[altedge] (kv) -- (k4) node[midway,above,font=\tiny] {1.0};

  \end{tikzpicture}
  \caption{K-mer graph for a simple SNV.  Reference edges (grey, weight 0.01)
    form the lower path; the variant k-mer $K_v$ (red, weight 1.0) forms an
    alternative path.  After reference-edge removal, the path
    BigBang $\to K_1 \to K_2 \to K_v \to K_4 \to$ BigCrunch is the only
    remaining path and represents the variant allele.}
  \label{fig:graph-example}
\end{figure}

\subsection{Variant Classification}
\label{sec:classification}

Each alternative path is compared to the reference path using the
\texttt{diff\_path\_without\_overlap} algorithm~\cite{audoux2020km}.  Let $R$
be the ordered reference k-mer sequence and $A$ the alternative path:

\begin{enumerate}
  \item Walk from the left to find the first divergence index $i$: the smallest
        index such that $R[i] \neq A[i]$.
  \item Walk from the right to find the last divergence indices $j_R$
        (reference) and $j_A$ (alternative).
  \item The DNA sequence of the variant region is reconstructed from the k-mers
        $A[i..j_A]$ by taking the first base of each k-mer and appending all
        bases of the last k-mer.
\end{enumerate}

Classification rules (Table~\ref{tab:classification}):

\begin{table}[h]
  \centering
  \caption{Variant classification rules in \textsc{diff\_path\_without\_overlap}.}
  \label{tab:classification}
  \begin{tabular}{lll}
    \toprule
    Condition & \kmerdet{} type & \km{} type \\
    \midrule
    $i = j_R = j_A$ & \texttt{Reference} & Reference \\
    $j_R = j_A$ (same length) & \texttt{Substitution} & Substitution \\
    $i = j_R^{\text{overlap}}$ & \texttt{Itd} & ITD \\
    $j_R < j_A$, no deleted bases & \texttt{Insertion} & Insertion \\
    $j_R > j_A$, no inserted bases & \texttt{Deletion} & Deletion \\
    Otherwise & \texttt{Indel} & Indel \\
    \bottomrule
  \end{tabular}
\end{table}

ITD detection relies on the overlap condition: if the divergence from the left
reaches the end of the reference k-mer sequence before the paths re-converge,
the variant path contains a duplication of the entire reference region, the
hallmark of an internal tandem duplication.

\subsubsection{INDEL Normalization}

INDELs produced by k-mer walking may not be in their canonical left-aligned
representation, particularly in homopolymer and short-repeat contexts.
\kmerdet{} applies the left-alignment normalization procedure of Tan et
al.~\cite{tan2015unified} to all insertion and deletion calls, shifting the
variant position as far left as possible while keeping the reference and
alternative alleles consistent.  This ensures compatibility with VCF conventions
and prevents duplicate calls for the same INDEL expressed differently.

\subsection{Quantification: NNLS Decomposition and rVAF}
\label{sec:nnls}

The \km{} algorithm models k-mer counts as a linear mixture of contributions
from all paths.  Let $m$ be the number of k-mer positions, $n$ be the number of
paths (reference + alternatives), and define:
\begin{itemize}
  \item $\mathbf{c} \in \mathbb{R}^m$: observed k-mer counts from $\mathcal{D}$.
  \item $C \in \{0,1,2,\ldots\}^{m \times n}$: contribution matrix where
    $C_{ij}$ counts how many times k-mer $i$ appears on path $j$.
  \item $\boldsymbol{\alpha} \in \mathbb{R}_{\geq 0}^n$: path coverage coefficients.
\end{itemize}

The model is $\mathbf{c} \approx C\boldsymbol{\alpha}$, and we seek the
non-negative solution minimising the residual:
\begin{equation}
  \hat{\boldsymbol{\alpha}} = \arg\min_{\boldsymbol{\alpha} \geq 0}
    \|C\boldsymbol{\alpha} - \mathbf{c}\|_2^2
  \label{eq:nnls}
\end{equation}

This is the standard NNLS problem, solved with the Lawson-Hanson
algorithm~\cite{lawson1974solving}.  The relative variant allele frequency is:
\begin{equation}
  \hat{f}_j = \frac{\hat{\alpha}_j}{\sum_{l} \hat{\alpha}_l}
  \qquad \forall j \in \{1,\ldots,n\}
  \label{eq:rvaf}
\end{equation}

Note that $C_{ij} > 1$ is possible for ITDs, where the duplicated segment
contributes k-mers to the variant path multiple times.  This is handled
naturally by the NNLS formulation.

\subsubsection{Bootstrap Confidence Intervals}

The point estimate $\hat{f}_j$ has no associated uncertainty.  \kmerdet{}
computes 95\% bootstrap confidence intervals by:
\begin{enumerate}
  \item Draw $B = 1000$ bootstrap resamples of the count vector $\mathbf{c}^*$
        by sampling independently from $\text{Poisson}(\hat{c}_i)$ for each
        position $i$.
  \item Solve NNLS for each resample to obtain $\hat{f}_j^{(b)}$.
  \item Report $[F_{2.5\%}, F_{97.5\%}]$ as the confidence interval.
\end{enumerate}

For a panel of 50 targets with $\sim$3 paths each, the $50 \times 3 \times 1000 =
150{,}000$ NNLS solutions (each for a system of size $\sim 31 \times 3$) complete
in $<1$ second with \texttt{ndarray-linalg} Lapack bindings.

\subsection{Statistical Scoring}
\label{sec:scoring}

For each variant call, \kmerdet{} computes:

\textbf{Per-k-mer binomial p-value}: At each variant k-mer position $i$ with
count $c_i$ and total sibling count $D_i$, compute $p_i$ from
equation~\eqref{eq:binomial} using $\varepsilon_k = \varepsilon_{\text{base}} / 3$.

\textbf{Combined p-value}: Fisher's method (equation~\eqref{eq:fisher}) over
the $m_{\text{eff}}$ approximately independent k-mer positions.

\textbf{Phred-scaled QUAL}: $Q = \min(999, -10 \log_{10} p_{\text{combined}})$.

\textbf{VCF FILTER field}: PASS if $Q \geq 30$ and min\_coverage $\geq 3$;
otherwise annotated with the failing criterion (LowQual, LowCoverage, etc.).

\subsection{Filtering Pipeline}
\label{sec:filtering}

\kmerdet{}'s \texttt{filter} subcommand implements a multi-stage pipeline:

\begin{enumerate}
  \item \textbf{Hard thresholds}: Minimum count (adaptive or fixed), minimum
    rVAF, minimum expression.
  \item \textbf{Type exclusion}: Remove \texttt{Reference} type calls; optionally
    restrict to specific variant types.
  \item \textbf{Context filters}: Elevated thresholds near homopolymer runs
    ($\geq 6$~bp), in low-complexity (linguistic complexity $< 0.8$) regions.
  \item \textbf{Panel-of-normals}: Remove calls present in $>$5\% of PoN
    samples.
  \item \textbf{Tumor-informed matching}: Reference mode (match by
    CHROM/POS/REF/ALT/TYPE) or alt-sequence mode (match by full variant sequence).
\end{enumerate}

Each stage annotates the call with its FILTER status, enabling retrospective
analysis of filter performance.

\subsection{Multi-k Detection}
\label{sec:multik}

To address the fundamental k-mer length constraint on INDEL sensitivity,
\kmerdet{} supports three multi-k strategies:

\subsubsection{Adaptive k Per Target}

Target metadata specifies a \texttt{recommended\_k} value based on variant type:
$k = 21$ for large insertions ($L > 15$~bp), $k = 31$ as the standard, $k = 41$
for targets in segmental duplications or high-repeat regions.  Each target
queries only one jellyfish database.

\subsubsection{Independent Multi-k with Union Merging}

Run the full pipeline at $k \in \{21, 31, 41\}$; variants detected at multiple
k values receive a detection-bonus multiplier (20\% per additional k).  Variant
matching across k values uses left-aligned genomic coordinates after INDEL
normalisation.

\subsubsection{Consensus Voting}

Rather than a simple union, consensus voting assigns tier labels based on the
number of k values detecting the variant (Table~\ref{tab:multikvoting}).

\begin{table}[h]
  \centering
  \caption{Multi-k consensus voting tiers.}
  \label{tab:multikvoting}
  \begin{tabular}{lll}
    \toprule
    Detection pattern & Confidence tier & Action \\
    \midrule
    All 3 k values & Tier~1 (highest) & PASS automatically \\
    2 of 3 k values & Tier~2 & PASS with note \\
    $k=31$ only & Tier~3 & PASS (standard) \\
    $k=21$ only & Tier~4 & Flag for review \\
    $k=41$ only & Tier~4 & Flag for review \\
    \bottomrule
  \end{tabular}
\end{table}

\subsection{Output Formats}
\label{sec:output}

\kmerdet{} produces output in five formats:

\begin{description}
  \item[TSV] Tab-separated values with the same column order as \km{}'s output
    (sample, target, type, variant\_name, rVAF, expression, min\_coverage, \ldots),
    extended with statistical annotation columns (pvalue, qual, ci\_lower,
    ci\_upper).  Fully backward-compatible with \km{}/\kmtools{} downstream tools.
  \item[CSV] Comma-separated equivalent of the TSV format.
  \item[VCF~4.3] Standard VCF with CHROM, POS, REF, ALT, QUAL, FILTER, INFO,
    and FORMAT fields.  INFO fields: DP (total depth), AF (rVAF), TLOD (tumour log
    odds).  FORMAT fields: AD (allelic depths), AF, CI (confidence interval bounds).
  \item[JSON/JSONL] Machine-readable format for programmatic consumption.
    JSONL (one record per line) is suitable for streaming and log-based
    monitoring systems.
  \item[XLSX] Excel workbook with one worksheet per sample, formatted tables, and
    per-column colour coding for rVAF values.
\end{description}
